\textbf{COLUNI 2026 -- Questão 14}\par\medskip Considere a classificação \textit{tri-perfeito} para qualquer número natural, diferente de zero, que satisfaça às seguintes características matemáticas:
\begin{enumerate}
\item[I.] É um quadrado perfeito, ou seja, pode ser escrito na forma $n^2$, $n\in\mathbb{N}^*$.
\item[II.] É o quadrado de uma soma finita de números naturais consecutivos, sendo o número 1 a primeira parcela dessa soma, ou seja, pode ser escrito na forma $(1+2+3+\cdots+n)^2$, $n\in\mathbb{N}^*$.
\item[III.] É igual a uma soma finita de cubos de números naturais consecutivos, sendo o número 1 a primeira parcela dessa soma, ou seja, pode ser escrito na forma $1^3+2^3+\cdots+n^3$, $n\in\mathbb{N}^*$.
\end{enumerate}
O ano de 2025 está sendo considerado como \textit{tri-perfeito}, pois satisfaz a todas as características matemáticas acima.

Considerando as informações fornecidas, é CORRETO afirmar que:
\begin{enumerate}
\item o ano de 2025 também possui a característica de ser um cubo perfeito.
\item o ano de 1826 foi classificado como \textit{tri-perfeito}.
\item o ano de 1946 apresentou uma das três características para ser \textit{tri-perfeito}.
\item o ano de 3025 será considerado \textit{tri-perfeito}.
\end{enumerate}
